\section{Global Smooth Solutions}
\label{app:smooth}

In this appendix we prove Theorem~\ref{thm:smooth}. The argument is a Picard
iteration through linear transport: each iterate is transported by the
velocity field of the previous one, the iterates obey bounds that are uniform
in the iteration index on every bounded time interval, and they converge
together with all their spatial derivatives. We write the iteration in
$L^2$-Sobolev form, so that, besides elementary real analysis, it uses only
Section~I.\ref{I-sec:prelim}, Section~\ref{sec:II-prelim} and the facts on ordinary differential equations
collected in Lemma~\ref{lem:ode}. For
the classical theory of smooth solutions see \cite{Wolibner}, \cite{Holder} and
\cite{Kato1967}.

\subsection{Ordinary Differential Equations}

The following facts are used without proof.

\begin{lemma}[Theorems 2.2, 2.10, 2.13 and 2.17 in \cite{Teschl}]
\label{lem:ode}
Let $n\ge1$ and $k\ge1$ be integers, let $U\subset\R\times\R^n$ be open, and let
$F\in C^k(U;\R^n)$.
\begin{enumerate}[label=\textup{(\roman*)}]
\item For every $(s,y)\in U$ the problem $\dot x=F(t,x)$, $x(s)=y$, has a unique
  maximal solution, defined on an open interval containing $s$. We write
  $\flow_F(t,s,y)$ for its value at time $t$.
\item Around each point of $U$ there is an open set $I\times\mathcal U\subset U$,
  with $I$ an interval, such that $\flow_F$ and $\partial_t\flow_F$ are defined
  and of class $C^k$ on $I\times I\times\mathcal U$, and there $D_y\flow_F$
  solves the variational equation
  \[
  \partial_tD_y\flow_F(t,s,y)=D_xF\big(t,\flow_F(t,s,y)\big)\,D_y\flow_F(t,s,y).
  \]
\item If $U=\R\times\R^n$ and for every $T>0$ there are numbers $M(T)$, $L(T)$
  with $|F(t,x)|\le M(T)+L(T)|x|$ on $[-T,T]\times\R^n$, then every maximal
  solution is defined on all of $\R$.
\end{enumerate}
\end{lemma}

Item (i) is Theorem 2.2 of \cite{Teschl} (the Picard-Lindel\"of theorem; a
$C^1$ function is locally Lipschitz) together with Theorem 2.13 there, which
produces the maximal solution. Item (ii) is Theorem 2.10 of \cite{Teschl}, and
item (iii) is Theorem 2.17 there. We use these facts through the following
lemma.

\begin{lemma}
\label{lem:flow}
Let $u\in C^\infty(\R\times\R^2;\R^2)$ be bounded on $[-T,T]\times\R^2$ for every
$T>0$, and write $\flow_u$ for the solution map of $\dot x=u(t,x)$ from
Lemma~\ref{lem:ode}(i). Then:
\begin{enumerate}[label=\textup{(\alph*)}]
\item $\flow_u(t,s,y)$ is defined for all $(t,s,y)\in\R\times\R\times\R^2$, and
  $\flow_u\in C^\infty(\R\times\R\times\R^2)$;
\item $\flow_u(t,\tau,\flow_u(\tau,s,y))=\flow_u(t,s,y)$ for all $t,\tau,s,y$, and
  each map $\flow_u(t,s,\cdot)$ is a $C^\infty$ diffeomorphism of $\R^2$ with
  inverse $\flow_u(s,t,\cdot)$;
\item $|\flow_u(t,s,y)-y|\le|t-s|\sup_{[-T,T]\times\R^2}|u|$ whenever
  $s,t\in[-T,T]$;
\item $Y(t)=D_y\flow_u(t,s,y)$ solves $\partial_tY=\nabla u(t,\flow_u(t,s,y))\,Y$
  with $Y(s)=I$, and if $\divop u=0$, then $\det D_y\flow_u\equiv1$.
\end{enumerate}
\end{lemma}

\begin{proof}
Lemma~\ref{lem:ode}(iii), with $M(T)=\sup_{[-T,T]\times\R^2}|u|$ and $L(T)=0$,
shows that $\flow_u$ is defined on $\R\times\R\times\R^2$. Both sides of the
identity in (b) solve $\dot x=u(t,x)$ as functions of $t$ and agree at $t=\tau$,
so they agree by the uniqueness in Lemma~\ref{lem:ode}(i). With $t=s$ the
identity gives $\flow_u(s,\tau,\flow_u(\tau,s,y))=y$, so $\flow_u(\tau,s,\cdot)$ is
a bijection with inverse $\flow_u(s,\tau,\cdot)$. Part (c) follows from
$\flow_u(t,s,y)-y=\int_s^tu(\sigma,\flow_u(\sigma,s,y))\dd\sigma$.

We prove the smoothness in (a). Fix $k\ge1$ and $(t_*,s_*,y_*)$, let $J$ be the
closed interval with endpoints $s_*$ and $t_*$, and put
$\gamma(\tau)=\flow_u(\tau,s_*,y_*)$. For every $\tau\in J$,
Lemma~\ref{lem:ode}(ii) gives an open set
$I_\tau\times\mathcal U_\tau\ni(\tau,\gamma(\tau))$ on which $\flow_u$ is $C^k$ in
the sense stated there. The sets
$\{\sigma\in J\cap I_\tau:\gamma(\sigma)\in\mathcal U_\tau\}$, $\tau\in J$, are
relatively open in $J$ and cover $J$. By the Lebesgue number lemma there are
points $s_*=\tau_0,\tau_1,\dots,\tau_N=t_*$, monotone in the index, and points
$\tau'_i\in J$ such that $\tau_i,\tau_{i+1}\in I_{\tau'_i}$ and
$\gamma(\tau_i)\in\mathcal U_{\tau'_i}$ for $0\le i<N$. By (b), for $(t,s,y)$ near $(t_*,s_*,y_*)$,
\[
\flow_u(t,s,y)=\flow_u\Big(t,\tau_{N-1},\flow_u\big(\tau_{N-1},\tau_{N-2},\cdots
\flow_u(\tau_2,\tau_1,\flow_u(\tau_1,s,y))\cdots\big)\Big).
\]
The innermost map $(s,y)\mapsto\flow_u(\tau_1,s,y)$ is $C^k$ near
$(\tau_0,\gamma(\tau_0))$; each map $y\mapsto\flow_u(\tau_{i+1},\tau_i,y)$ is $C^k$
on $\mathcal U_{\tau'_i}$, which is a neighborhood of $\gamma(\tau_i)$, and it
maps $\gamma(\tau_i)$ to $\gamma(\tau_{i+1})$; and
$(t,x)\mapsto\flow_u(t,\tau_{N-1},x)$ is $C^k$ for $t\in I_{\tau'_{N-1}}\ni t_*$
and $x\in\mathcal U_{\tau'_{N-1}}$. By
continuity and the chain rule, $\flow_u$ is $C^k$ near $(t_*,s_*,y_*)$. Since $k$
is arbitrary, $\flow_u\in C^\infty$, and by (b) the maps $\flow_u(t,s,\cdot)$ are
$C^\infty$ diffeomorphisms.

We prove (d). By (a), $\flow_u$ is $C^\infty$. Differentiating
$\partial_t\flow_u(t,s,y)=u(t,\flow_u(t,s,y))$ in $y$ and exchanging the order
of the derivatives gives
\[
\partial_tD_y\flow_u(t,s,y)=\nabla u\big(t,\flow_u(t,s,y)\big)\,D_y\flow_u(t,s,y),
\]
and $Y(s)=I$ because $\flow_u(s,s,y)=y$. The matrix
$Y(t)$ is invertible by (b), and Jacobi's formula gives
$\partial_t\det Y=\det Y\operatorname{tr}(Y^{-1}\partial_tY)=\det
Y\operatorname{tr}\nabla u=(\divop u)\det Y$. If $\divop u=0$, then $\det Y$ is
constant in $t$ and equal to $\det Y(s)=1$.
\end{proof}

\subsection{Two Estimates on the Plane}

\begin{lemma}
\label{lem:grad-bound}
Let $g\in C^1_c(\R^2)$ be real with $g\not\equiv0$, and let $u=\KBS\ast g$. Then
for all $x\in\R^2$ and $i,l\in\{1,2\}$,
\[
|\partial_lu_i(x)|\le\|g\|_{L^\infty}\Big(2+\log_+\frac{\|\nabla g\|_{L^\infty}}{\|g\|_{L^\infty}}\Big)
+\frac{\|g\|_{L^1}}{2\pi},
\]
where $\|\nabla g\|_{L^\infty}=\sup_x|\nabla g(x)|$.
\end{lemma}

\begin{proof}
Since $\KBS$ is locally integrable and $g\in C^1_c$, differentiation under the
integral sign gives $\partial_lu_i(x)=\int(\KBS)_i(z)\,\partial_lg(x-z)\dd z$.
Fix $\delta\in(0,1]$. By Lemma~I.\ref{I-lem:BS-basic}(a), the part of the integral
over $\{|z|<\delta\}$ is at most $\delta\|\nabla g\|_{L^\infty}$. On
$\{|z|>\delta\}$ we write $\partial_lg(x-z)=-\partial_{z_l}[g(x-z)]$ and apply
the divergence theorem on $\{\delta<|z|<\rho'\}$, with $\rho'$ so large that
$g(x-z)=0$ for $|z|\ge\rho'$. The outward normal on $\{|z|=\delta\}$ is
$-z/|z|$, so
\begin{align*}
\int_{|z|>\delta}(\KBS)_i(z)\,\partial_lg(x-z)\dd z
&=\int_{|z|>\delta}\partial_l(\KBS)_i(z)\,g(x-z)\dd z\\
&\quad+\int_{|z|=\delta}(\KBS)_i(z)\,g(x-z)\frac{z_l}{|z|}\dd S .
\end{align*}
By Lemma~I.\ref{I-lem:BS-basic}(a), the boundary term is at most
$(2\pi\delta)^{-1}\|g\|_{L^\infty}2\pi\delta=\|g\|_{L^\infty}$, the volume integral
over $\{\delta<|z|<1\}$ is at most
$\|g\|_{L^\infty}\int_\delta^1r^{-1}\dd r=\|g\|_{L^\infty}\log\frac1\delta$, and the
volume integral over $\{|z|\ge1\}$ is at most $\|g\|_{L^1}/(2\pi)$. Hence
$|\partial_lu_i(x)|\le\delta\|\nabla g\|_{L^\infty}+\|g\|_{L^\infty}\big(1+\log\frac1\delta\big)+\|g\|_{L^1}/(2\pi)$.
The choice $\delta=\min\{1,\|g\|_{L^\infty}/\|\nabla g\|_{L^\infty}\}$ gives the
claim.
\end{proof}

\begin{lemma}
\label{lem:comm-R2}
Let $s\ge2$ be an integer, $g\in C_c^\infty(\R^2)$, and let
$u\in C^\infty(\R^2;\R^2)$ satisfy
$|\nabla u|_{\max}:=\max_{i,l}\sup_x|\partial_lu_i(x)|<\infty$ and
$\partial^\beta u\in L^2$ for $1\le|\beta|\le s$. Then for every multi-index
$\alpha$ with $|\alpha|=s$,
\[
\|[\partial^\alpha,u\cdot\nabla]g\|_{L^2}\le2^{s+1}\Big[\big(|\nabla u|_{\max}+\|\nabla u\|_{H^{s-1}}\big)\|g\|_{H^s}
+\|\nabla u\|_{H^{s-1}}\|\nabla g\|_{L^\infty}\Big].
\]
\end{lemma}

\begin{proof}
By the Leibniz formula,
$[\partial^\alpha,u\cdot\nabla]g=\sum_{l=1}^2\sum_{0<\beta\le\alpha}\binom\alpha\beta\partial^\beta u_l\,\partial_l\partial^{\alpha-\beta}g$,
and $\sum_l\sum_{\beta\le\alpha}\binom\alpha\beta=2^{s+1}$. It suffices to bound
each product by the bracket on the right. If $|\beta|=1$, the product is at
most $|\nabla u|_{\max}\|g\|_{H^s}$. If $\beta=\alpha$, it is at most
$\|\partial^\alpha u_l\|_{L^2}\|\partial_lg\|_{L^\infty}\le\|\nabla u\|_{H^{s-1}}\|\nabla g\|_{L^\infty}$,
since $\partial^\alpha u_l$ is a derivative of order $s-1$ of an entry of
$\nabla u$. If $2\le|\beta|\le s-1$, H\"older's inequality and (E1) of
Lemma~\ref{lem:embed} give the bound
$\|\partial^\beta u_l\|_{H^1}\|\partial_l\partial^{\alpha-\beta}g\|_{H^1}
\le\|\nabla u\|_{H^{|\beta|}}\|g\|_{H^{s-|\beta|+2}}\le\|\nabla u\|_{H^{s-1}}\|g\|_{H^s}$,
because $|\beta|\le s-1$ and $|\beta|\ge2$.
\end{proof}

\subsection{The Iteration}

Throughout the rest of this appendix, $\omega_0$ and $R_0$ are as in
Theorem~\ref{thm:smooth}, and we assume $\omega_0\not\equiv0$. We put
\begin{gather*}
A=\|\omega_0\|_{L^\infty},\qquad
N_0=\sup_x|\nabla\omega_0(x)|,\\
U_*=\Big(\frac{2A\|\omega_0\|_{L^1}}{\pi}\Big)^{1/2},\qquad
c_*=2+\frac{\|\omega_0\|_{L^1}}{2\pi A},
\end{gather*}
and, for $t\in\R$,
\begin{gather*}
\mathcal L(t)=\Big(\log_+\frac{N_0}{A}+c_*\Big)e^{2A|t|}-c_*,\qquad
\Psi(t)=Ae^{\mathcal L(t)},\\
V(t)=2A\big(\mathcal L(t)+c_*\big),\qquad R(t)=R_0+U_*|t| .
\end{gather*}
The functions $\mathcal L$, $\Psi$ and $V$ are even and nondecreasing in
$|t|$; moreover $\mathcal L\ge\mathcal L(0)=\log_+(N_0/A)\ge0$, and
$\mathcal L'=V$ on $(0,\infty)$. Since $\supp\omega_0$ is a compact subset of
$B_{R_0}$, there is
$R_0'<R_0$ with $\omega_0=0$ outside $\overline B_{R_0'}$.

\begin{definition}
\label{def:iterates}
Put $\omega^{[0]}(t,x)=\omega_0(x)$ for $(t,x)\in\R\times\R^2$. If
$\omega^{[j]}\in C^\infty(\R\times\R^2)$ is defined and
$v^{[j]}=\KBS\ast\omega^{[j]}$ satisfies the hypotheses of Lemma~\ref{lem:flow},
let $\flow^{[j]}=\flow_{v^{[j]}}$ and put
$\omega^{[j+1]}(t,x)=\omega_0\big(\flow^{[j]}(0,t,x)\big)$.
\end{definition}

\begin{lemma}
\label{lem:iterates}
The iterates $\omega^{[j]}$ are defined for every $j\ge0$, and:
\begin{enumerate}[label=\textup{(\roman*)}]
\item $\omega^{[j]}\in C^\infty(\R\times\R^2)$ is real, and $\omega^{[j]}(t,x)=0$
  for $|x|>R_0'+U_*|t|$;
\item $\|\omega^{[j]}(t)\|_{L^p}=\|\omega_0\|_{L^p}$ for every $t\in\R$ and every
  $p\in[1,\infty]$;
\item $v^{[j]}\in C^\infty(\R\times\R^2;\R^2)$, every derivative of $v^{[j]}$ is
  bounded on $[-T,T]\times\R^2$ for every $T>0$, $\divop v^{[j]}=0$, and
  $|v^{[j]}|\le U_*$;
\item $\flow^{[j]}$ is defined and $C^\infty$ on $\R\times\R\times\R^2$,
  $\det D_y\flow^{[j]}\equiv1$, $\omega^{[j+1]}(0)=\omega_0$, and
  $\partial_t\omega^{[j+1]}+v^{[j]}\cdot\nabla\omega^{[j+1]}=0$ on $\R\times\R^2$.
\end{enumerate}
\end{lemma}

\begin{proof}
We argue by induction on $j$. Items (i) and (ii) hold for $j=0$. Assume (i) and
(ii) for some $j$.

We prove (iii) for $j$. Fix $T>0$. Every derivative
$\partial_t^i\partial_x^\gamma\omega^{[j]}$ is continuous and vanishes outside
the compact set $[-T,T]\times\overline B_{R(T)}$ on $[-T,T]\times\R^2$, so it is
bounded there. When $(t,x)$ ranges over a bounded subset of
$(-T,T)\times\R^2$, the functions
$z\mapsto\KBS(z)\,\partial_t^i\partial_x^\gamma\omega^{[j]}(t,x-z)$ are dominated
by a constant multiple of the integrable function
$|\KBS(z)|\mathbf 1_{\{|z|\le\rho\}}$, for a suitable $\rho$. Differentiation under the integral sign
and dominated convergence give
$\partial_t^i\partial_x^\gamma v^{[j]}=\KBS\ast\partial_t^i\partial_x^\gamma\omega^{[j]}$,
continuous, and bounded on $[-T,T]\times\R^2$ by Lemma~I.\ref{I-lem:BS-basic}(b).
Lemma~I.\ref{I-lem:BS-basic}(d) gives $\divop v^{[j]}=0$, and
Lemma~I.\ref{I-lem:BS-basic}(b) with (ii) gives
\[
|v^{[j]}(t,x)|\le\big(2\|\omega^{[j]}(t)\|_{L^1}\|\omega^{[j]}(t)\|_{L^\infty}/\pi\big)^{1/2}=U_* .
\]

We prove (iv) for $j$. By (iii), Lemma~\ref{lem:flow} applies to $v^{[j]}$ and
gives the properties of $\flow^{[j]}$. Hence $\omega^{[j+1]}$ is defined, real
and $C^\infty$, and $\omega^{[j+1]}(0,x)=\omega_0(x)$. By Lemma~\ref{lem:flow}(b),
$\omega^{[j+1]}(t,\flow^{[j]}(t,0,a))=\omega_0(a)$ for all $t$ and $a$.
Differentiating in $t$ gives
$(\partial_t\omega^{[j+1]}+v^{[j]}\cdot\nabla\omega^{[j+1]})(t,\flow^{[j]}(t,0,a))=0$,
and $\flow^{[j]}(t,0,\cdot)$ is onto.

We prove (i) and (ii) for $j+1$. If $\omega^{[j+1]}(t,x)\ne0$, then
$a=\flow^{[j]}(0,t,x)$ satisfies $|a|\le R_0'$, and by Lemma~\ref{lem:flow}(c),
$|x-a|\le U_*|t|$; so $|x|\le R_0'+U_*|t|$. The substitution
$x=\flow^{[j]}(t,0,a)$, whose Jacobian determinant is $1$, gives
$\int|\omega^{[j+1]}(t,x)|^p\dd x=\int|\omega_0(a)|^p\dd a$ for $p<\infty$, and the
case $p=\infty$ follows because $\flow^{[j]}(0,t,\cdot)$ is a bijection.
\end{proof}

\begin{lemma}
\label{lem:iter-lip}
For every $j\ge0$, $t\in\R$ and $x\in\R^2$,
\[
|\nabla\omega^{[j]}(t,x)|\le\Psi(t),\qquad
\max_{i,l}|\partial_lv^{[j]}_i(t,x)|\le A\big(\mathcal L(t)+c_*\big),\qquad
\|\nabla v^{[j]}(t,x)\|_{\mathrm{op}}\le V(t).
\]
\end{lemma}

\begin{proof}
The third bound follows from the second, because the operator norm of a
$2\times2$ matrix is at most twice the largest modulus of its entries. We show,
by induction on $j$, that the first bound holds for all $t$ and that it implies
the second.

For $j=0$ we have
$|\nabla\omega_0|\le N_0\le\max\{A,N_0\}=Ae^{\mathcal L(0)}\le\Psi(t)$. Suppose
that $|\nabla\omega^{[j]}(\tau,x)|\le\Psi(\tau)$ for all $\tau$ and $x$. Since
$\mathcal L\ge0$, we have $\log_+(\Psi(\tau)/A)=\mathcal L(\tau)$. By
Lemma~\ref{lem:iterates}(ii), $\|\omega^{[j]}(\tau)\|_{L^\infty}=A$ and
$\|\omega^{[j]}(\tau)\|_{L^1}=\|\omega_0\|_{L^1}$. Thus Lemma~\ref{lem:grad-bound},
applied to $g=\omega^{[j]}(\tau)$, gives
\[
|\partial_lv^{[j]}_i(\tau,x)|\le A\big(2+\mathcal L(\tau)\big)+\frac{\|\omega_0\|_{L^1}}{2\pi}=A\big(\mathcal L(\tau)+c_*\big).
\]
Hence $\|\nabla v^{[j]}(\tau,x)\|_{\mathrm{op}}\le V(\tau)$.

Next, by the chain rule,
$\nabla\omega^{[j+1]}(t,x)=D_y\flow^{[j]}(0,t,x)^{\mathsf T}\,\nabla\omega_0\big(\flow^{[j]}(0,t,x)\big)$.
By Lemma~\ref{lem:flow}(d), $Y(\tau)=D_y\flow^{[j]}(\tau,t,x)$ solves
$\partial_\tau Y=\nabla v^{[j]}(\tau,\flow^{[j]}(\tau,t,x))Y$ with $Y(t)=I$, so
$\|Y(\tau)\|_{\mathrm{op}}\le1+\big|\int_t^\tau V(\sigma)\|Y(\sigma)\|_{\mathrm{op}}\dd\sigma\big|$,
and Gronwall's inequality gives
$\|Y(0)\|_{\mathrm{op}}\le\exp\big|\int_0^tV(\sigma)\dd\sigma\big|=e^{\mathcal L(t)-\mathcal L(0)}$,
because $\mathcal L'=V$ on $(0,\infty)$ and both functions are even. Hence
$|\nabla\omega^{[j+1]}(t,x)|\le N_0e^{\mathcal L(t)-\mathcal L(0)}\le Ae^{\mathcal L(t)}=\Psi(t)$,
since $N_0e^{-\mathcal L(0)}=\min\{N_0,A\}$.
\end{proof}

\begin{lemma}
\label{lem:iter-Hs}
For every integer $s\ge1$ and every $T>0$, and for all $j\ge0$ and $|t|\le T$,
$\|\omega^{[j]}(t)\|_{H^s}\le\mathcal M_s(T)$, where
\begin{gather*}
\mathcal M_1(T)^2=\|\omega_0\|_{L^2}^2+\pi R(T)^2\Psi(T)^2,\\
\mathcal M_s(T)^2=\big(\|\omega_0\|_{H^s}^2+\mathcal M_{s-1}(T)^2+1\big)e^{\mathcal B_s(T)T}+\mathcal M_{s-1}(T)^2\quad(s\ge2),\\
\mathcal B_s(T)=(s+1)2^{s+2}\big(V(T)+\mathcal M_{s-1}(T)+\mathcal M_{s-1}(T)\Psi(T)\big).
\end{gather*}
In particular $\mathcal M_s(T)$ depends only on $s$, $T$, $R_0$, $A$, $N_0$ and
the norms $\|\omega_0\|_{L^1}$ and $\|\omega_0\|_{H^{s'}}$, $0\le s'\le s$.
\end{lemma}

\begin{proof}
We argue by induction on $s$. For $s=1$, Lemma~\ref{lem:iterates}(i)--(ii) and
Lemma~\ref{lem:iter-lip} give $\|\omega^{[j]}(t)\|_{L^2}=\|\omega_0\|_{L^2}$ and
$\|\nabla\omega^{[j]}(t)\|_{L^2}^2\le|B_{R(T)}|\Psi(T)^2$.

Let $s\ge2$ and assume the bound at level $s-1$ for all iterates. For $j=0$
the bound holds because $\mathcal M_s(T)\ge\|\omega_0\|_{H^s}$. Let $j\ge0$ and
$|\alpha|=s$. Differentiating the equation of Lemma~\ref{lem:iterates}(iv) gives
\[
\partial_t\partial^\alpha\omega^{[j+1]}+v^{[j]}\cdot\nabla\partial^\alpha\omega^{[j+1]}
=-[\partial^\alpha,v^{[j]}\cdot\nabla]\omega^{[j+1]} .
\]
We multiply by $\partial^\alpha\omega^{[j+1]}$ and integrate; all functions are
smooth and real, and $\omega^{[j+1]}(t)$ vanishes outside $B_{R(T)}$ for
$|t|\le T$. Since $\divop v^{[j]}=0$, the transport term gives
$-\frac12\int(\divop v^{[j]})(\partial^\alpha\omega^{[j+1]})^2\dd x=0$. Put
$\mathcal H(t)=\sum_{|\alpha|=s}\|\partial^\alpha\omega^{[j+1]}(t)\|_{L^2}^2$ and
$Z(t)=\|\omega^{[j+1]}(t)\|_{H^s}$; there are $s+1$ multi-indices with
$|\alpha|=s$. Put $\mathcal M=\mathcal M_{s-1}(T)$. We apply Lemma~\ref{lem:comm-R2} with
$u=v^{[j]}(t)$ and $g=\omega^{[j+1]}(t)$, using $|\nabla v^{[j]}|_{\max}\le V(T)$
and $|\nabla\omega^{[j+1]}|\le\Psi(T)$ (Lemma~\ref{lem:iter-lip}), and
\[
\|\nabla v^{[j]}(t)\|_{H^{s-1}}=\|\omega^{[j]}(t)\|_{H^{s-1}}\le\mathcal M
\]
(Lemma~I.\ref{I-lem:BS-basic}(d) and the induction hypothesis). This gives, for
$|t|\le T$,
\begin{align*}
|\mathcal H'(t)|&\le(s+1)2^{s+2}\big[(V(T)+\mathcal M)Z+\mathcal M\Psi(T)\big]Z\\
&\le\mathcal B_s(T)(Z^2+1)\le\mathcal B_s(T)\big(\mathcal H+\mathcal M^2+1\big),
\end{align*}
because $Z\le Z^2+1$ and $Z^2\le\mathcal H+\|\omega^{[j+1]}(t)\|_{H^{s-1}}^2\le\mathcal H+\mathcal M^2$.
Since $\mathcal H(0)\le\|\omega_0\|_{H^s}^2$, Gronwall's inequality gives
$\mathcal H(t)\le(\|\omega_0\|_{H^s}^2+\mathcal M^2+1)e^{\mathcal B_s(T)|t|}$, and
$Z^2\le\mathcal H+\mathcal M^2\le\mathcal M_s(T)^2$.
\end{proof}

\subsection{Convergence}

\begin{lemma}
\label{lem:iter-cauchy}
Let $T>0$ and put $e_j(t)=\|\omega^{[j]}(t)-\omega^{[j-1]}(t)\|_{L^2}$ for $j\ge1$.
Then for every $j\ge1$ and every $t\in[-T,T]$,
\[
e_{j}(t)\le2\|\omega_0\|_{L^2}\frac{\big(2R(T)\Psi(T)|t|\big)^{j-1}}{(j-1)!} .
\]
\end{lemma}

\begin{proof}
Put $R=R(T)$ and $\Psi=\Psi(T)$. For $j=1$, $e_1(t)\le2\|\omega_0\|_{L^2}$ by
Lemma~\ref{lem:iterates}(ii). Let $j\ge1$. Subtracting the equations of
Lemma~\ref{lem:iterates}(iv) for $\omega^{[j+1]}$ and $\omega^{[j]}$ shows that
$\delta=\omega^{[j+1]}-\omega^{[j]}$ satisfies
\[
\partial_t\delta+v^{[j]}\cdot\nabla\delta=-\big(\KBS\ast(\omega^{[j]}-\omega^{[j-1]})\big)\cdot\nabla\omega^{[j]} .
\]
We multiply by $\delta$ and integrate; the transport term vanishes, because
$\divop v^{[j]}=0$ and $\delta(t)$ has compact support. For $|t|\le T$, $\nabla\omega^{[j]}(t)$ vanishes outside $B_R$ and is
bounded by $\Psi$, and Lemma~I.\ref{I-lem:BS-basic}(c), with $\rho=R$, gives
$\|\KBS\ast(\omega^{[j]}-\omega^{[j-1]})(t)\|_{L^2(B_R)}\le2Re_j(t)$. Hence
$|\frac{\dd}{\dd t}e_{j+1}^2|\le4R\Psi e_je_{j+1}$, and for all $\epsilon>0$,
$\big|\frac{\dd}{\dd t}(e_{j+1}^2+\epsilon)^{1/2}\big|\le2R\Psi e_j$. Since
$e_{j+1}(0)=0$, letting $\epsilon\to0$ gives
$e_{j+1}(t)\le2R\Psi\big|\int_0^te_j(\tau)\dd\tau\big|$ for $|t|\le T$, and the
claim follows by induction on $j$.
\end{proof}

\begin{proposition}
\label{prop:iter-limit}
There is a real $\omega\in C^\infty(\R\times\R^2)$ such that, with
$v=\KBS\ast\omega$, for every $T>0$ and every multi-index $\gamma$,
$\partial_x^\gamma\omega^{[j]}\to\partial_x^\gamma\omega$ and
$\partial_x^\gamma v^{[j]}\to\partial_x^\gamma v$ uniformly on $[-T,T]\times\R^2$
as $j\to\infty$. The function $\omega$ has properties \textup{(1)}--\textup{(3)}
of Theorem~\ref{thm:smooth}, and for all $t\in\R$, $x\in\R^2$ and
$p\in[1,\infty]$,
\begin{equation}
\label{eq:limit-bounds}
\begin{gathered}
\omega(t,x)=0\ \text{ for }|x|>R_0'+U_*|t|,\qquad
\|\omega(t)\|_{L^p}=\|\omega_0\|_{L^p},\qquad |\nabla\omega(t,x)|\le\Psi(t),\\
|v(t,x)|\le U_*,\qquad \max_{i,l}|\partial_lv_i(t,x)|\le A\big(\mathcal L(t)+c_*\big).
\end{gathered}
\end{equation}
\end{proposition}

\begin{proof}
We first prove the convergence. Fix $T>0$. By Lemma~\ref{lem:iter-cauchy} the
series $\sum_j\sup_{|t|\le T}e_j(t)$ is at most
$2\|\omega_0\|_{L^2}\exp(2R(T)\Psi(T)T)$. Thus $(\omega^{[j]})$ is a Cauchy sequence
for the norm $\sup_{|t|\le T}\|\cdot\|_{L^2}$. For an integer
$k\ge1$, (E3) of Lemma~\ref{lem:embed} and Lemma~\ref{lem:iter-Hs} give
\[
\sup_{|t|\le T}\|\omega^{[j]}(t)-\omega^{[i]}(t)\|_{H^k}^2\le3^k\cdot2\mathcal M_{2k}(T)\sup_{|t|\le T}\|\omega^{[j]}(t)-\omega^{[i]}(t)\|_{L^2},
\]
which tends to $0$ as $i,j\to\infty$. By (E2), for all $\gamma$ the functions
$\partial_x^\gamma\omega^{[j]}$, which are continuous on $\R\times\R^2$, form a
uniform Cauchy sequence on $[-T,T]\times\R^2$. Their uniform limits are
continuous; the limit of $\partial_x^\gamma\omega^{[j]}$ is $\partial_x^\gamma$ of
the limit $\omega$ of $\omega^{[j]}$, as one sees by passing to the limit in
$g(x+he_l)-g(x)=\int_0^h\partial_lg(x+\sigma e_l)\dd\sigma$, with
$g=\partial_x^\gamma\omega^{[j]}(t)$; and the limits for
different $T$ agree. The first line of Equation~\eqref{eq:limit-bounds} follows
from Lemma~\ref{lem:iterates}(i)--(ii) and Lemma~\ref{lem:iter-lip}, because the
convergence is uniform and the supports lie in a fixed ball for $|t|\le T$. In
particular $\omega(t)\in C_c^\infty(\R^2)$ for all $t$, and $\omega$ has
property (1) of Theorem~\ref{thm:smooth}, since $R_0'+U_*|t|<R(t)$.

We now treat the velocity. For a bounded measurable $g$ vanishing outside
$B_\rho$, Lemma~I.\ref{I-lem:BS-basic}(b) and $\|g\|_{L^1}\le\pi\rho^2\|g\|_{L^\infty}$
give $\|\KBS\ast g\|_{L^\infty}\le\sqrt2\,\rho\|g\|_{L^\infty}$. Applied to
$g=\partial_x^\gamma\omega^{[j]}(t)-\partial_x^\gamma\omega(t)$ with $\rho=R(T)$,
this shows that $\partial_x^\gamma v^{[j]}=\KBS\ast\partial_x^\gamma\omega^{[j]}$
converges to $\KBS\ast\partial_x^\gamma\omega$ uniformly on $[-T,T]\times\R^2$.
By Lemma~I.\ref{I-lem:BS-basic}(d) applied to $\omega(t)$, $\KBS\ast\partial_x^\gamma\omega(t)=\partial_x^\gamma v(t)$
and $\divop v=0$. The bounds on $v$ in Equation~\eqref{eq:limit-bounds} follow
from Lemma~\ref{lem:iterates}(iii) and Lemma~\ref{lem:iter-lip}.

We now derive the equation. For every $\gamma$, Lemma~\ref{lem:iterates}(iv)
gives
\[
\partial_x^\gamma\omega^{[j+1]}(t,x)=\partial^\gamma\omega_0(x)
-\int_0^t\partial_x^\gamma\big(v^{[j]}\cdot\nabla\omega^{[j+1]}\big)(\tau,x)\dd\tau .
\]
By the Leibniz formula the integrand is a finite sum of products of
$x$-derivatives of $v^{[j]}$ and of $\omega^{[j+1]}$, each converging uniformly
and bounded on $[-T,T]\times\R^2$. Passing to the limit,
\[
\partial_x^\gamma\omega(t,x)=\partial^\gamma\omega_0(x)-\int_0^t\partial_x^\gamma(v\cdot\nabla\omega)(\tau,x)\dd\tau,
\]
with an integrand that is continuous and bounded on $[-T,T]\times\R^2$. Hence
$\partial_t\partial_x^\gamma\omega$ exists, is continuous and bounded, and
equals $-\partial_x^\gamma(v\cdot\nabla\omega)$. With $\gamma=0$ this is
property (3).

Finally we prove that $\omega$ and $v$ are $C^\infty$ with bounded derivatives.
For $T>0$ and an integer $l\ge0$, let $\mathcal C^l(T)$ be the set of functions
$g$ on $(-T,T)\times\R^2$ such that $\partial_t^i\partial_x^\gamma g$ exists, is
continuous and is bounded, for all $i\le l$ and all multi-indices $\gamma$
(including $i=0$ and $\gamma=0$); here the
$x$-derivatives are taken first. We have shown $\omega\in\mathcal C^1(T)$. Suppose
$\omega\in\mathcal C^l(T)$. Every $\partial_t^i\partial_x^\gamma\omega$ with
$i\le l$ vanishes for $|x|\ge R(T)$, so differentiation under the integral sign,
as in the proof of Lemma~\ref{lem:iterates}(iii), gives
$\partial_t^i\partial_x^\gamma v=\KBS\ast\partial_t^i\partial_x^\gamma\omega$,
continuous and bounded by $\sqrt2R(T)\sup|\partial_t^i\partial_x^\gamma\omega|$;
so $v\in\mathcal C^l(T)$. By the Leibniz formula $v\cdot\nabla\omega\in\mathcal C^l(T)$,
and the identity $\partial_t\partial_x^\gamma\omega=-\partial_x^\gamma(v\cdot\nabla\omega)$
shows that $\partial_t^{i+1}\partial_x^\gamma\omega$ exists, is continuous and
is bounded for $i\le l$. Hence $\omega\in\mathcal C^{l+1}(T)$. By induction,
$\omega$ and $v$ lie in $\mathcal C^l(T)$ for all $l$ and $T$. By Schwarz's
theorem on mixed partial derivatives (if $\partial_ag$, $\partial_bg$ and
$\partial_b\partial_ag$ exist near a point and $\partial_b\partial_ag$ is
continuous at it, then $\partial_a\partial_bg$ exists at that point and equals
$\partial_b\partial_ag$ there), an induction on the order shows that every
partial derivative of $\omega$ or $v$, taken in any order, exists, is continuous,
and equals the corresponding $\partial_t^i\partial_x^\gamma$ derivative. Thus
$\omega,v\in C^\infty(\R\times\R^2)$, $v$ and every derivative of $v$ are bounded
on $[-T,T]\times\R^2$ for all $T$, and $\omega$ has property (2).
\end{proof}

\subsection{Proof of Theorem~\ref{thm:smooth}}

The proof of Theorem~\ref{thm:smooth} also gives the following explicit bounds.

\begin{lemma}
\label{lem:smooth-bounds}
In Theorem~\ref{thm:smooth}, assume $\omega_0\not\equiv0$ and put
$A=\|\omega_0\|_{L^\infty}$, $N_0=\sup_x|\nabla\omega_0(x)|$,
$U_*=\sqrt{2A\|\omega_0\|_{L^1}/\pi}$, $c_*=2+\|\omega_0\|_{L^1}/(2\pi A)$,
$\mathcal L(t)=(\log_+(N_0/A)+c_*)e^{2A|t|}-c_*$ and $\Psi(t)=Ae^{\mathcal L(t)}$.
Then for every $t\in\R$ and $x\in\R^2$,
\begin{gather*}
|\nabla\omega(t,x)|\le\Psi(t),\qquad|v(t,x)|\le U_*,\\
\max_{i,l}|\partial_lv_i(t,x)|\le A\big(\mathcal L(t)+c_*\big),\qquad
\|\nabla v(t,x)\|_{\mathrm{op}}\le2A\big(\mathcal L(t)+c_*\big).
\end{gather*}
\end{lemma}

\begin{proof}[Proof of Theorem~\ref{thm:smooth} and Lemma~\ref{lem:smooth-bounds}]
We first prove (5), for any $\omega$ with properties (1)--(3). Let $\tilde\omega$
be as in (5), and fix a compact interval $J\subset I$ containing $0$. Let
$\rho>0$ be such that $\omega(t)$ and $\tilde\omega(t)$ vanish outside $B_\rho$ for
$t\in J$; this is possible by (1) and the hypothesis on $\tilde\omega$. Put
$\delta=\omega-\tilde\omega\in C^1(J\times\R^2)$. Since
$\partial_t\tilde\omega=-v\cdot\nabla\tilde\omega+(v-\KBS\ast\tilde\omega)\cdot\nabla\tilde\omega$,
\[
\partial_t\delta+v\cdot\nabla\delta=-(\KBS\ast\delta)\cdot\nabla\tilde\omega
\qquad\text{on }J\times\R^2 .
\]
We multiply by $\delta$ and integrate in $x$. The transport term vanishes,
because $\divop v=0$ and $\delta(t)$ has compact support. The function
$\nabla\tilde\omega$ is continuous and vanishes outside $J\times B_\rho$, so it is
bounded on $J\times\R^2$ by some $L<\infty$, and
Lemma~I.\ref{I-lem:BS-basic}(c) gives $\|\KBS\ast\delta(t)\|_{L^2(B_\rho)}\le2\rho\|\delta(t)\|_{L^2}$.
Hence $|\frac{\dd}{\dd t}\|\delta(t)\|_{L^2}^2|\le4\rho L\|\delta(t)\|_{L^2}^2$ on
$J$. Since $\delta(0)=0$, Gronwall's inequality gives $\delta=0$ on $J$. Every
point of $I$ lies in such an interval $J$, which proves (5). In particular at
most one $\omega\in C^\infty(\R\times\R^2)$ has properties (1)--(3), since such
an $\omega$ satisfies the hypotheses on $\tilde\omega$ in (5) with $I=\R$.

If $\omega_0\equiv0$, the function $\omega\equiv0$ has properties (1)--(4), (6)
and (7). From now on $\omega_0\not\equiv0$, and $\omega$ is the function of
Proposition~\ref{prop:iter-limit}, which has properties (1)--(3). The bounds of
Lemma~\ref{lem:smooth-bounds} are contained in Equation~\eqref{eq:limit-bounds},
together with $\|M\|_{\mathrm{op}}\le2\max_{i,l}|M_{il}|$ for $2\times2$ matrices
$M$.

We prove (4). By (2), Lemma~\ref{lem:flow} applies to $v$; let
$\flow(t,y)=\flow_v(t,0,y)$ and $\flow_t=\flow(t,\cdot)$. By Lemma~\ref{lem:flow},
$\flow$ is defined for all $t$; by (a) and (b) there, $\flow$ and
$(t,x)\mapsto\flow_t^{-1}(x)=\flow_v(0,t,x)$ lie in $C^\infty(\R\times\R^2)$ and
each $\flow_t$ is a $C^\infty$ diffeomorphism; and by (d),
$\det D\flow_t\equiv1$ because $\divop v=0$. By (3),
$\frac{\dd}{\dd t}\omega(t,\flow(t,a))=(\partial_t\omega+v\cdot\nabla\omega)(t,\flow(t,a))=0$,
so $\omega(t)=\omega_0\circ\flow_t^{-1}$. For a continuous $F$ with $F(0)=0$,
the function $F(\omega(t))$ is continuous with compact support, and the change
of variables $x=\flow_t(a)$ gives $\int F(\omega(t))\dd x=\int F(\omega_0)\dd a$.

We prove (6). Let $\Orot_\theta\in SO(2)$. Since $(\Orot_\theta z)^\perp=\Orot_\theta z^\perp$,
we have $\KBS(\Orot_\theta z)=\Orot_\theta\KBS(z)$. Put
$\omega_\theta(t,x)=\omega(t,\Orot_\theta^{-1}x)$. The substitution
$y=\Orot_\theta y'$ gives
$\KBS\ast\omega_\theta(t,x)=\Orot_\theta v(t,\Orot_\theta^{-1}x)$, and
$\nabla\omega_\theta(t,x)=\Orot_\theta\nabla\omega(t,\Orot_\theta^{-1}x)$. Hence
$(\KBS\ast\omega_\theta)\cdot\nabla\omega_\theta=(v\cdot\nabla\omega)(t,\Orot_\theta^{-1}x)$,
and $\omega_\theta$ is a $C^\infty$ solution of the equation in (3) with datum
$\omega_0\circ\Orot_\theta^{-1}$, vanishing outside $B_{R(t)}$ at time $t$. If
$\omega_0\circ\Orot_\theta=\omega_0$, then $\omega_0\circ\Orot_\theta^{-1}=\omega_0$,
and (5) with $I=\R$ gives $\omega_\theta=\omega$, that is,
$\omega(t)\circ\Orot_\theta^{-1}=\omega(t)$ for all $t$, which is equivalent to
$\omega(t)\circ\Orot_\theta=\omega(t)$.

We prove (7). Let $T>0$, $\varphi\in\DT$, and let $\Ren\in C^1(\R)$ with
$\Ren'$ bounded and $\Ren(0)=0$. The function $\Ren(\omega)$ lies in
$C^1(\R\times\R^2)$, vanishes outside $B_{R(t)}$ at time $t$, and by (3)
satisfies $\partial_t\Ren(\omega)+v\cdot\nabla\Ren(\omega)=\Ren'(\omega)(\partial_t\omega+v\cdot\nabla\omega)=0$.
We multiply by $\varphi$ and integrate over $(0,T)\times\R^2$. Since
$\varphi(t,\cdot)=0$ for $t$ near $T$, an integration by parts in $t$ gives
$\int_0^T\!\int\partial_t\Ren(\omega)\,\varphi=-\int\Ren(\omega_0)\varphi(0,\cdot)-\int_0^T\!\int\Ren(\omega)\,\partial_t\varphi$,
and since $\divop v=0$, an integration by parts in $x$ gives
$\int(v\cdot\nabla\Ren(\omega))\varphi\dd x=-\int\Ren(\omega)\,v\cdot\nabla\varphi\dd x$.
Hence Equation~(I.\ref{I-eq:renorm}) holds, and with $\Ren(s)=s$ so does
Equation~(I.\ref{I-eq:weak}). Finally, $\|\omega(t)\|_{L^1}=\|\omega_0\|_{L^1}$ by
(4) with $F(s)=|s|$, and $\|\omega(t)\|_{L^\infty}=\|\omega_0\|_{L^\infty}$
because $\omega(t)=\omega_0\circ\flow_t^{-1}$. Moreover $\omega$ is continuous and
$\omega(t)$ vanishes outside $B_{R(T)}$ for $|t|\le T$, so $\omega$ is uniformly
continuous on $[0,T]\times\R^2$, and $t\mapsto\omega(t)$ is continuous from
$[0,T]$ into $(L^1\cap L^q)(\R^2)$ for all $q\in[1,\infty]$. Hence the
restriction of $\omega$ to $[0,T)$ lies in $L^\infty_t(L^1\cap L^\infty)_x$, and
in $X_q(T)$ for every $q\in(4/3,\infty)$. This proves (7) and completes the
proof.
\end{proof}
